%% problem_id: S048-PJM-305-1-p01-lemma-12
%% source_id: S048 (MSP; Pacific Journal of Mathematics)
%% source_article: Kenneth L. Baker, "The Poincare homology sphere, lens space surgeries, and some
%%   knots with tunnel number two", Pacific J. Math. 305 (2020), no. 1, 1--30.
%% source_url: https://msp.org/pjm/2020/305-1/pjm-v305-n1-p01-s.pdf
%% representation: PDF; transcribed by rendering the page (pdftoppm -r 150 -png) and reading the
%%   image, then checking against it.
%% statement_locator: PDF p. 10 (printed p. 9), Lemma 12
%% proof_locator: PDF p. 10 (printed p. 9), the Proof immediately following
%% difficulty_level: H2
%% FIGURE DEPENDENCE: the proof cannot be read without Figure 3 (printed p. 9).  Per plan 4.3 the
%%   figure must be preserved, so the page image is kept as raw/msp-p10.png and must be read with this
%%   item.  This is why the item is not presented as self-contained.
%% status: complete (statement + author proof verbatim + reason + locators + figure pointer)
%% ---------------------------------------------------------------------------
\documentclass[11pt]{article}
\usepackage{amsmath,amssymb,amsthm}
\newtheorem{lemma}{Lemma}
\begin{document}

\section*{Lemma 12}

\textbf{Source.} K.~L.~Baker, \emph{The Poincar\'e homology sphere, lens space surgeries, and some
knots with tunnel number two}, Pacific J.\ Math.\ \textbf{305} (2020), no.~1; printed page 9.

\begin{lemma}
The knot $K_n^*$ in $L(3n^2+n+1,\,-3n+2)$ that is surgery dual to $K_n$ is homologous to the knot
$T_n^*$ that is surgery dual to the $(3n+1,n)$-torus knot in $S^3$.
\end{lemma}

\subsection*{Proof (authors' text)}

Starting from the lower left of Figure 1 where the dual arc $\kappa_n^*$ is presented on a standard
form of the two bridge link $B(3n^2+n+1,\,-3n+2)$, Figure 3 (left) shows how two clasp changes (and
isotopy) transforms $\kappa_n^*$ into an arc $\tau_n^*$ in a bridge sphere. Since $\tau_n^*$ lies in
the bridge sphere, its lift to the double branched cover is a knot $T_n^*$ in the Heegaard torus of
the lens space. Lifting the transformation of Figure 3 (left) to the double branched cover thus shows
that $K_n^*$ and $T_n^*$ are related by two crossing changes. (The clasp changes lift to crossing
changes as indicated in Figure 3 (right).) Hence these knots are homotopic, and therefore homologous,
in the lens space. \hfill$\square$

\medskip
\noindent\textit{Note.} Figure 3 lives on the same printed page; see \texttt{raw/msp-p10.png}.

\end{document}
